\section*{\Large Appendix}
\vspace{0.1in}

\section{Some Ancillary Material}
\label{sec:ancillary}

\subsection{Review of \gpucb}
\label{sec:gpucbReview}

We present a review of the \gpucbs algorithm of~\citet{srinivas10gpbandits} which we build
on in this work.
Here we will assume $\func\sim\GP(\zero,\kernel)$ where $\kernel:\Xcal^2\rightarrow\RR$ is
a radial kernel defined on the domain $\Xcal$.
The algorithm is given below.

\insertAlgorithmGPBALGO

To present the theoretical results for \gpucb, we begin
by defining the \emph{Maximum Information Gain} (\mig) which characterises the
statistical difficulty of GP bandits.

\begin{definition}(Maximum Information Gain~\citep{srinivas10gpbandits})
Let $f\sim \GP(\zero, \phix)$.
Consider any $A\subset\RR^d$ and
let $A' = \{x_1, \dots, x_n\} \subset A$ be a finite subset.
Let $f_{A'}, \epsilon_{A'}\in\RR^n$ such that $(f_{A'})_i=f(x_i)$ and
$(\epsilon_{A'})_i\sim\Ncal(0,\eta^2)$.
Let $y_{A'} = f_{A'}+\epsilon_{A'}$.
Denote the Shannon Mutual Information by $I$.
The Maximum Information Gain of $A$ is
\[
\IG_n(A) = \max_{A'\subset A, |A'| = n} I(y_{A'}; f_{A'}).
\]
\label{def:infGain}
\end{definition}
\insertpostspacing

Next, we will need the following regularity conditions on the kernel.
It is satisfied for four times differentiable
kernels such as the SE kernel and \matern{} kernel when
$\nu>2$~\citep{ghosal06gpconsistency}.

% \insertprespacing
\begin{assumption} Let $\func\sim\GP(\zero,\kernel)$, where
$\kernel:\Xcal^2 \rightarrow \RR$ is a stationary kernel.
%  $\kernel(\cdot, x)$ is $L$-Lipschitz for all $x$. 
The partial derivatives
of $\func$  satisfies the following condition.
There exist constants $a, b >0$ such that,
\[
\text{for all $J>0$, $\;$and for all $i \in \{1,\dots,d\}$},\quad \PP\left( \sup_{x} 
\Big|\partialfrac{x_i}{\func(x)}\Big| > J \right)
\leq a e^{-(J/b)^2}.
\]
\label{asm:kernelAssumption}
\end{assumption}

The following theorem is a bound on the simple regret $S_n$~\eqref{eqn:SnDefn} for \gpucb.

\begin{theorem}(\cite{srinivas10gpbandits})
\label{thm:gpucb}
Let $\func\sim\GP(\zero,\kernel)$, where $\Xcal=[0,1]^d$,
$\func:\Xcal\rightarrow\RR$ and the kernel $\kernel$ satisfies
Assumption~\ref{asm:kernelAssumption}).
At each query, we have noisy observations
$y = f(x) + \epsilon$ where $\epsilon\sim\Ncal(0,\eta^2)$.
Denote $C_1 = 8/\log(1+\eta^{-2})$.
Pick a failure probability $\delta\in(0,1)$ and run \emph{\gpucbs} with
  $\betat = 2\log\left(\frac{2\pi^2t^2}{3\delta}\right) + 
  2d\log\left(t^2bdr\sqrt{\frac{4ad}{\delta}}\right)$.
The following holds with probability $>1-\delta$,
% \hspace{-0.2in}
\vspace{-0.02in}
\[
% \PP\left( \forall n\geq 1,\; R_n \leq \sqrt{C_1n\betan\IGn(\Xcal)} \,+\, 2 \right)
% \geq 1-\delta
\text{for all $n\geq 1$,}
\hspace{0.4in}
S_n \leq \sqrt{\frac{C_1\betan\IGn(\Xcal)}{n}} \,+\, \frac{\pi^2}{6}.
\]
\end{theorem}


\subsection{Some Technical Results}
\label{sec:technical}

Here we present some technical lemmas we will need for our analysis.

\begin{lemma}[Gaussian Concentration]
\label{lem:gaussConcentration}
Let $Z\sim\Ncal(0,1)$. Then $\;\PP(Z>\epsilon) \leq \frac{1}{2}\exp(-\epsilon^2/2)$.
\end{lemma}
% \insertprespacing

\begin{lemma}[Mutual Information in GP, \cite{srinivas10gpbandits} Lemma 5.3]
\label{lem:IGformula}
Let $\func\sim\GP(\zero,\kernel)$, $f:\Xcal\rightarrow\RR$ and we observe
$y=f(x) + \epsilon$ where $\epsilon\sim\Ncal(0,\eta^2)$. Let $A$ be a finite
subset of $\Xcal$ and $f_A,y_A$ be the function values and observations on this set
respectively. Then the Shannon Mutual Information
$I(y_A;f_A)$ is,
\[
I(y_A; \func_A) = \frac{1}{2} \sum_{t=1}^{n} \log(1 + \eta^{-2}
  \sigma^2_{t-1}(x_t) ).
\]
where $\sigma^2_{t-1}$ is the posterior GP variance after observing the first $t-1$
points. 
\end{lemma}

Our next result is a technical lemma taken from~\citet{kandasamy16mfbo}.
It will be used in controlling
the posterior variance of our $\func$ and $\gunc$ GPs.

\begin{lemma}[Posterior Variance Bound~\citep{kandasamy16mfbo}]
\label{lem:varbound}
Let $\func\sim(\zero,\kernel)$, $\func:\Ucal\rightarrow\RR$ where
$\kernel(u,u') = \kernelscale\phi(\|u-u'\|)$ and
$\phi$ is a radial kernel.
Upon evaluating $f$ at $u$ we observe $y = f(u) + \epsilon$ where
$\epsilon\sim\Ncal(0,\eta^2)$.
Let $u_1\in\Ucal$ and suppose we have $s$ observations at $u_1$ and no observations
elsewhere.
Then the posterior variance $\kernel'$ (see~\eqref{eqn:gpPost}) at all $u\in\Ucal$
satisfies,
\[
\kernel'(u,u) \leq \kernelscale(1-\phi^2(\|u-u_1\|)) + \frac{\eta^2/s}{ 1+
\frac{\eta^2}{\kernelscale s}}.
\]
\end{lemma}
\textbf{Proof:}
The proof is in Section C.0.1 of~\citet{kandasamy16mfbo} who
prove this result as part of a larger proof.
\hfill\BlackBox


